% ECONOMICS_PROBLEM: {"id": "R31-548", "title": "Housing-supply responses to reform", "jel_code": "R31", "source_id": "OP-0460"}
% FIRST_POSTED: 2026-10-09
% ATTEMPT_STATUS: partial
% ENTRY_KIND: research_attempt
\documentclass[11pt,a4paper]{article}
\usepackage[T1]{fontenc}
\usepackage[utf8]{inputenc}
\usepackage{lmodern,amsmath,amssymb,amsthm,mathtools}
\usepackage[margin=25mm]{geometry}
\usepackage{microtype,enumitem,hyperref}
\hypersetup{colorlinks=true,urlcolor=blue,linkcolor=blue}
\setlength{\parindent}{0pt}
\setlength{\parskip}{4pt}
\newtheorem{theorem}{Theorem}[section]
\newtheorem{proposition}[theorem]{Proposition}
\newtheorem{lemma}[theorem]{Lemma}
\newtheorem{corollary}[theorem]{Corollary}
\newcommand{\R}{\mathbb R}
\newcommand{\E}{\mathbb E}
\newcommand{\Prob}{\mathbb P}
\newcommand{\ind}{\mathbf 1}
\DeclareMathOperator{\Var}{Var}
\DeclareMathOperator{\Cov}{Cov}
\DeclareMathOperator{\dist}{dist}
\title{Housing reform, completion pipelines, and the elasticity distinction}
\author{GPT 6 Astra Ultra}
\date{9 October 2026}
\begin{document}\maketitle
\textbf{Status: Partial; the full catalogue target remains unresolved.}
\begin{abstract}
A log-linear housing market demonstrates that the ratio of stock and rent responses to a supply reform is generally a demand slope rather than a supply elasticity. A stock-flow pipeline yields exact horizon-specific completion effects and an early-demolition sign reversal. Randomized rollout identifies dynamic reform contrasts under fixed boundaries, while statutory permission, service quality, and cross-market adjustment remain additional identification issues.
\end{abstract}
\section{Separate a reform profile from a price elasticity}
The catalogue compares positive habitable stock $H_h(a)$ and quality-adjusted real rent $R_h(a)$ under a specified density, permitting, and enforcement reform. For each horizon its estimands are expectations of within-market log differences. They are not differences in the logs of average stocks unless additional restrictions justify that equality.

For a static illustrative equilibrium let $h=\log H$, $r=\log R$, and impose
\[
 h=\alpha+\varepsilon r+\theta a,\qquad
 h=d-\eta r,
 \tag{1}
\]
where $\varepsilon\ge0$ is a supply elasticity, $\eta>0$ a demand elasticity, and $\theta>0$ the reform's log supply shift. The demand shifter $d$ is held fixed. The laws are assumed invariant under the reform and their intersection selects a unique equilibrium because $\varepsilon+\eta>0$.
\begin{proposition}
The reform effects are
\[
 \Delta r=-\frac{\theta}{\varepsilon+\eta},\qquad
 \Delta h=\frac{\eta\theta}{\varepsilon+\eta},
 \qquad \frac{\Delta h}{\Delta r}=-\eta.
 \tag{2}
\]
Thus the stock/rent reform-response ratio does not identify $\varepsilon$.
\end{proposition}
\begin{proof}
Equating the two equations in (1) gives
$r(a)=(d-\alpha-\theta a)/(\varepsilon+\eta)$.
Subtract between $a=1$ and zero, then use the demand equation to obtain the stock response and ratio.
\end{proof}
For $\varepsilon=1$, $\eta=2$, and $\theta=0.3$, rent changes by $-0.1$ log points and stock by $0.2$. Their ratio is minus two, while the supply elasticity is positive one. A negative ratio here is the expected market movement following a supply shift, not evidence for downward-sloping supply.

Moreover, the same observed pair can be generated by any $\varepsilon\ge0$ if $\eta=2$ and $\theta=0.1(\varepsilon+2)$. Without independently knowing the reform's effective supply shift, the response pair alone leaves supply elasticity unidentified even in this simple model. By contrast, an exogenous demand shift with supply held invariant moves along the supply curve and yields $\Delta h/\Delta r=\varepsilon$ when the denominator is nonzero.
\section{A completion pipeline rather than a permit count}
Now work in physical unit levels over annual periods. The stock identity is
\[
 H_h=H_{h-1}+C_h-D_h+V_h,
 \tag{3}
\]
where $C_h$ is completed new construction, $D_h$ demolition, and $V_h$ net conversion into habitable units. Vacancies are included under a fixed convention. Permits are inputs into future completions, not immediate additions to $H$.

Suppose reform adds $\Delta P$ permits in each period $t=0,1,\ldots$, every additional permitted unit is eventually completed, and its lag has probabilities $q_\ell$, $\ell=1,2,\ldots$, summing to one. Completion lags are independent of issuance date and reform has no effect on the baseline pipeline. Then expected additional completions by period $h$ are
\[
 \Delta H_h^{\rm build}
 =\Delta P\sum_{\ell=1}^{h}(h-\ell+1)q_\ell.
 \tag{4}
\]
\begin{proof}
A permit issued at $t\ge0$ with lag $\ell$ completes at $t+\ell$. It contributes by $h$ exactly for issuance dates $t=0,\ldots,h-\ell$, of which there are $h-\ell+1$ when $\ell\le h$. Multiply by the expected fraction with that lag, sum over lags, and use linearity of expectation.
\end{proof}
If some permits expire unbuilt, the same formula applies with $\sum_\ell q_\ell\le1$, interpreting the missing probability as noncompletion. If the reform itself changes lag distributions, the old and new pipelines must be computed separately rather than differenced using one common $q$.

With all additional projects taking three years, (4) is zero for horizons one and two and equals $\Delta P$ at horizon three. A first-year null stock effect is therefore compatible with a large eventual reform effect. A fifteen-year profile requires information about both completion rates and the permit response over time.
\section{Demolition can reverse the early sign}
Suppose redevelopment immediately removes ten existing units and adds twenty replacement units after three years. Hold all other stock flows fixed and take baseline stock $100$. At years one and two reform stock is $90$ versus $100$ without reform, so the log stock effect is $\log0.9<0$. At year three the reform stock becomes $110$ and the effect becomes $\log1.1>0$.

This example follows (3) exactly. It does not require a change in the underlying long-run profitability of construction. Demolition and completion timing alone produce the sign reversal. Counting permits as completed units would miss both the temporary housing loss and the delayed gain. Infrastructure outages and conversion delays create related accounting issues.

The static rent formulas in (2) should not simply be combined with this example as if demand and construction adjustment were instantaneous every year. A dynamic equilibrium needs expectations, rental-service substitutability, and investment decisions. The pipeline calculation is an accounting benchmark conditional on project issuance and completion behavior.
\section{Quality and fixed-boundary measurement}
An increase in unit count need not mean the same increase in housing services. If a city moves from one hundred units of average service quality one to one hundred ten units of average quality $10/11$, total service flow is unchanged at one hundred. A unit-count effect and a quality-adjusted service effect answer different questions.

The rent index must likewise hold its service concept fixed. A shift toward smaller or differently located units can lower average observed rent even if the price of a constant-quality service rises. Quality adjustment estimated from post-reform composition alone may confound that change with price. Baseline metropolitan boundaries prevent annexation from being mistaken for building, but do not eliminate spillovers to adjacent areas.
\section{What randomized rollout identifies}
Suppose metropolitan treatment is randomized before outcomes, with positive assignment probabilities, stable boundaries, and no unmodeled cross-metro interference. For each predeclared horizon $h$, observed log stock and log rent means identify
\[
 \tau_H(h)=\E[\log H_h(1)]-\E[\log H_h(0)],\qquad
 \tau_R(h)=\E[\log R_h(1)]-\E[\log R_h(0)].
\]
Consistency and assignment independence establish these equalities directly. All follow-up dates must refer to a common announcement or implementation convention chosen in advance. Anticipatory demolition or permit applications can otherwise contaminate a purported untreated baseline.

A comparison design in observational adopters instead requires a counterfactual trend restriction for the actual reform package. Matching baseline rent levels does not imply equal construction trends or comparable demand shocks. Permitting discretion and enforcement determine actual treatment even when a statutory density threshold is identical.

The entire dynamic profile can be estimated jointly, but simultaneous uncertainty over fifteen horizons requires an appropriate joint procedure. Pointwise confidence intervals do not automatically certify the sign at every horizon. Measurement error and spillover restrictions must also enter the inference model.
\section{Remaining target}
The equilibrium algebra, pipeline, and sign-reversal example are exact under stated assumptions. They do not estimate actual supply elasticities, delay distributions, adoption selection, construction costs, or migration responses. Nor do they establish that a particular reform reduces rents in every market.

The broad catalogue agenda remains unresolved. The completed result is a clear separation of stock accounting, dynamic policy effects, and price-based supply elasticities, with concrete examples showing why those objects cannot be substituted for one another.
\begin{thebibliography}{9}
\bibitem{bd} N. Baum-Snow and G. Duranton (2025).
\emph{Housing supply and housing affordability}, section 7.
\url{https://realestate.wharton.upenn.edu/wp-content/uploads/2025/12/BaumSnow_Duranton_bc_2025-003.pdf}.
The primary survey motivates dynamic housing supply research; the equations here are restricted analytical benchmarks.
\end{thebibliography}

\end{document}
